% Fragment: §5-§6 of Bounded Deontics. Preamble, macros, theorem envs live in bd-draft.tex.
% (For standalone compilation, prepend the bd-draft.tex preamble + \setcounter{section}{4}.)

\section{The derived necessity relation: bounded \textsc{must} as goal-domination}
\label{sec:nec}

The object level of a contract or statute, in our setting, is a \emph{labelled
transition system}---the ``letter'' of the law as a machine of states and admissible
actions, carrying no goals and no preferences of its own. Deontic force is then not
something we \emph{declare} on that machine but something we \emph{derive} from it,
relative to a goal supplied separately. This section makes that precise; the next
locates its boundary.

\begin{definition}[Object level]
A contract is a labelled transition system $T=(S,A,{\to},s_0)$ with states $S$, actions
$A$, transition relation ${\to}\subseteq S\times A\times S$, and initial state $s_0$. A
\emph{goal} $J\subseteq S$ is a set of states (equivalently, a state proposition). A path
$\pi$ from $s$ \emph{reaches} $J$ if some state of $\pi$ lies in $J$; an action $K$
\emph{occurs on} $\pi$ if some transition of $\pi$ is $K$-labelled. We write
$s\models\EF J$ when some path from $s$ reaches $J$.
\end{definition}

\begin{definition}[Goal-domination and the necessity relation]
\label{def:dom}
Action $K$ \emph{dominates} goal $J$ from state $s$, written $K\in\dom_s(J)$, iff every
path from $s$ that reaches $J$ contains a $K$-labelled transition. The \emph{necessity
relation} at $s$ is
\[
  \nect_s \;=\; \{(K,J)\in A\times \mathcal{G} : K\in\dom_s(J)\},
\]
over the family $\mathcal{G}$ of goals of interest. The two bounded modalities are read
off this relation:
\[
  s\models \MAY\,J \;\;\Longleftrightarrow\;\; s\models \EF J,
  \qquad
  s\models \MUST\,K \text{ for } J \;\;\Longleftrightarrow\;\; K\in\dom_s(J).
\]
That is: you \emph{may} pursue a goal when it is still reachable; you \emph{must} do
exactly those acts without which it is not.
\end{definition}

\paragraph{A rule-free instance.} The derived reading is starkest where no rule is present
at all. A restaurant patron may raise a hand, fold a napkin, or let a spoon fall---all
permitted, none prescribed; the object level is a bare menu of acts with no obligation
written into it. Fix a goal $J$, \emph{another glass of wine}, and the house convention that
wine is summoned from the waiter rather than fetched from the cellar. Then exactly one of
those permitted acts lies on every path to $J$, so that $\dom_s(J)=\{\textit{summon the
waiter}\}$---of which raising a hand is the conventional token---while the napkin and the
spoon, equally permitted, are not promoted, because nothing in \emph{their} relation to $J$
makes them necessary. This is the whole content of the claim that the \MUST{} is
\emph{derived} from the letter and a goal rather than \emph{declared}; and were the cellar
self-service, a second route to $J$ would open and the \MUST{} would relax to a \MAY{}.

\begin{remark}[Provenance: we concede the construct, and claim only its reading]
\label{rem:landmarks}
Definition~\ref{def:dom} is, as a graph construct, not new. An action on every path to
the goal is precisely a \emph{landmark} in automated planning \src{Hoffmann, Porteous \&
Sebastia 2004}, whose minimal-cut form is made explicit by \textsc{lm}-cut
\src{Helmert \& Domshlak 2009}; computing dominators is near-linear
\src{Lengauer \& Tarjan 1979}. We claim none of this machinery. Our contribution begins
only with the \emph{reading} of this party-neutral necessity as the deontic \MUST{}, and
with what that reading makes possible in Section~\ref{sec:boundary}. We flag one homonym
to avoid: our goal-\emph{dominator} is unrelated to the decision-theoretic
\emph{dominance} of one action over another across states of nature \src{Horty 2001}.
\end{remark}

\begin{remark}[Letter and spirit: the object level stays goal-free]
The goal $J$ does not appear in $T$. It enters only at the assertion level, as the
query against which $\nect$ is computed---the ``spirit'' read against the ``letter.''
This is the point of departure from accounts that attach the goal \emph{to} the
obligation, as in the dyadic instrumental obligation $O(J\!:\!K)$ of \src{Yan \& He 2025},
where the goal $J$ is an argument of the operator (a causal target the action must bring
about); here the obligation is \emph{derived}, the object level stays goal-free, and a
single contract answers arbitrarily many goal queries.
\end{remark}

\begin{remark}[The relation is many-to-many]
$\nect_s$ is a genuine relation, not a function. One action may dominate many goals; one
goal may have many dominators; and whether $K$ dominates $J$ is \emph{non-local}---it
turns on the presence or absence of \emph{alternative} paths to $J$ elsewhere in $T$,
not on anything local to $K$. (This non-locality is also why necessity does not compose
across conjunctions of goals; see Remark~\ref{rem:basket}.)
\end{remark}

\paragraph{Deadlines, and the non-temporal $\MAY\!\to\!\MUST$ transition.}
Because $\nect_s$ is indexed by the state $s$, the bounded modality of an action can
change as the system evolves. An action that is avoidable while alternative routes to
$J$ survive becomes \emph{necessary} at the state where the last such alternative lapses:
it \emph{enters} $\dom_s(J)$. A \textsc{within} deadline is one cause of this
$\MAY\!\to\!\MUST$ transition, but it is not the only one---the alternatives may instead
be eliminated by a change of circumstances rather than the passage of time. (Our running
example in Section~\ref{sec:boundary}, a protester whose only remaining route to a bed
for the night is arrest, is exactly such a \emph{non-temporal} collapse.) One caveat we
honour: the necessary set grows monotonically only for a \emph{cumulative, whole-trace}
notion of necessity; for the \emph{from-here-onward} $\nect_s$ defined above, an act
already performed drops out of future necessity, and the two must not be conflated.

\begin{remark}[Why not a lattice of goal-baskets]
\label{rem:basket}
It is tempting to lift $\nect$ to \emph{baskets} of goals---``what is the minimal set of
acts I am bound to perform to secure all of $J_1,\dots,J_n$?''---via a concept lattice
over $(A\times\mathcal{G},\nect)$. This is unsound: the lattice computes the
column-intersection $\bigcap_i \dom_s(J_i)$, which can strictly \emph{under}-report the
joint-goal necessity (take $a;c$ reaching $J_1$, $b;c$ reaching $J_2$, and only $a;b;c$
reaching both: the intersection yields $\{c\}$ where the correct bound is $\{a,b,c\}$).
Because necessity is non-local it does not factor through per-cell incidence. We
therefore present $\nect$ only at a fixed state-cut and defer the basket question---whose
correct tool is minimal-hitting-set enumeration \src{Reiter 1987}, not a Galois
lattice---to separate work.
\end{remark}

% The MAY dual of the dominator MUST, and how the prior art meets the cost of
% checking it (van der Meyden's EXPTIME DLP; MCMAS symbolic checking; planning's
% relaxation-based landmark extraction). Sits at the end of the necessity section
% because it is the existential companion to the dominator; forward-refers sec:tooling.
% =====================================================================
% SECTION DRAFT: "The permission dual, and the cost of checking it"
% Related-work paragraph on how the prior art meets state-space explosion,
% for the bounded-deontics paper. Slots into the background material near
% Section~\ref{sec:nec}/\ref{sec:tooling}, or as a subsection of the
% composition/tooling discussion. \src{Author Year} placeholders match the
% draft convention; bib entries (van der Meyden 1996/2005/2025) merged into
% ../bounded-deontics/draft/bounded-deontics.bib.
%
% ACCURACY NOTE: van der Meyden's DLP classifies EXECUTIONS (state-sequences),
% not states, as permitted, and distinguishes free-choice permission from
% permission-as-lack-of-prohibition. It is NOT goal-indexed. Do not conflate it
% with the Anderson/Meyer <a>not-V reading, and do not describe it as
% "the edge that keeps the goal reachable" -- that is our reconstruction, not
% his result.
% =====================================================================

\subsection{The permission dual, and the cost of checking it}
\label{sec:explosion}

If a \textsc{must} is the dominator of a goal---an act on \emph{every} accepting
run (Section~\ref{sec:nec})---then its dual is immediate and old. A \textsc{may}
is the existential companion: an act that lies on \emph{some} accepting run, an
edge whose taking does not foreclose the goal. Read the necessity/possibility
pair over a transition system and this is just $\Box/\Diamond$: the
Anderson--Kanger reduction already gives $O(A)=\Box(\lnot A\to V)$ with dual
$P(A)=\Diamond(A\land\lnot V)$ \src{Anderson 1958}, and Meyer's dynamic-logic
recasting makes the edge explicit, $F\alpha=[\alpha]V$ and
$P\alpha=\langle\alpha\rangle\lnot V$ \src{Meyer 1988}. The permission-as-goal-non-precluding reading we use is that
diamond, specialised to a goal automaton; we claim no novelty for it.

The sharpest transition-system treatment of permission is van der Meyden's
\emph{dynamic logic of permission} \src{van der Meyden 1996}, and it is worth
stating precisely, because it is both closer and more careful than the slogan.
Van der Meyden classifies not states but \emph{executions}---sequences of
states---as permitted or not, so that a permission constrains the intermediate
course of an action and not merely its result; and he separates
\emph{free-choice} permission from permission-as-lack-of-prohibition, the former
being exactly what standard modal logics cannot capture. His reading is
therefore trajectory-indexed rather than goal-indexed: it is the home of
free-choice permission, not of our goal-reachability dual, and the two should not
be run together. (The same Ron van der Meyden later reduces \textsc{dlp} to
\textsc{pdl} \src{van der Meyden 2005}, and, three decades on, is the co-author of
the executable-SAFE case study \src{van der Meyden \& Maher 2025} nearest our own
smart-contract work---one author spanning the classical semantics of permission
and its contemporary executable use.)

That precision comes at a price, and the price is the point. Van der Meyden's
\textsc{dlp} is \textsc{exptime}-complete, and his contribution is a complete
axiomatisation, not a checker that scales: the classical deontic tradition, from
Anderson through Meyer to van der Meyden, meets state-space explosion by
\emph{not enumerating a model at all}---it works at the level of proof theory and
Kripke semantics, where the exponential lives in the decision procedure rather
than in a graph one builds. Where the normative content \emph{is} model-checked,
in the normative multi-agent tradition---deontic interpreted systems and their
\textsc{mcmas}-style verifiers \src{Lomuscio \& Sergot 2003}---explosion is met
with the generic apparatus of symbolic model checking: \textsc{obdd} state
representation, partial-order reduction, bounded model checking. Nothing there is
specific to the deontic operators; the mitigation is the standard temporal-logic
toolbox pointed at a normative property.

The dominator reading buys something neither of those does, and it buys it by
borrowing. A dominator is a \emph{structural} property of the transition
graph, computable in near-linear time \src{Lengauer \& Tarjan 1979}; the
exponential never lived in the dominator computation but in \emph{building} the
graph, the product of the parties' state spaces. The one tradition that both
frames obligation as a necessary passage and has a native way to dodge that
product is AI planning, where our dominator is the \emph{landmark}---a fact true
on every plan \src{Hoffmann, Porteous \& Sebastia 2004}. Planning does not
compute landmarks on the true reachable graph; it extracts sound-but-incomplete
landmarks from the delete-relaxed problem in polynomial time
\src{Helmert \& Domshlak 2009}, and rides their count as a search heuristic
\src{Richter \& Westphal 2010}. The permission dual inherits the mirror
machinery: proving that an edge \emph{does} preclude the goal is dead-end
detection, itself relaxation- and abstraction-based. Reading \textsc{must} as
domination thus imports, as a side effect, planning's relaxation-based answer to
explosion---an answer the $\Box/\Diamond$ tradition could not reach because it
never made the landmark move. That import, not the duality, is what
Section~\ref{sec:tooling} turns into a checker.


\section{The boundary: two orderings, and where the reduction holds}
\label{sec:boundary}

The construction of Section~\ref{sec:nec} is, we freely concede, a \emph{synthesis} of
established parts rather than a new primitive. Our claim is the composition and its
jurisprudential reading. We therefore proceed by conceding each ingredient by name, and
then isolating the four places where the assembled object does something its parts,
taken singly, do not.

\subsection{The ordering source, indexed by party}
\label{sec:ordering}

Reachability fixes what is \emph{necessary} to attain a goal, but not \emph{which} goals
are worth attaining. For that we add a second component, in the manner of
\src{Kratzer 1991}: a \emph{modal base} (here the transition system, party-neutral) and
an \emph{ordering source}. We index the ordering source by party: each party $p$ carries
a preorder $\ppref$ over outcomes. Writing $\Goals_p(s)$ for the $\ppref$-best outcomes
reachable from $s$, the \emph{teleological} obligation on $p$ is
\[
  \MUST_p(s)\;=\;\bigcup_{J\in\Goals_p(s)} \dom_s(J).
\]
Both halves are anticipated and conceded: ordering sources indexed by an agent's goals
are familiar from anankastic conditionals \src{von Fintel \& Iatridou 2005;
Condoravdi \& Lauer 2016}, and the composition of an agent's ordering with a
necessity-to-ensure over a transition system is, for a single agent, essentially
\src{Shea-Blymyer \& Abbas 2022}. Our residue is narrow but real: the necessity we
compose is a \emph{graph dominator}---an act on \emph{every} path---rather than a
Bellman-optimal action, a best-world selection, or a causal intervention. It is this
dominator form that the constructions of Sections~\ref{sec:duality}--\ref{sec:force}
require, and that the neighbours lack. (A terminological caution: \src{Shea-Blymyer \&
Abbas 2022} build on Horty's \emph{Dominance} Act Utilitarianism, so their ``dominance''
is utility-dominance over actions---a second homonym, distinct from our graph
goal-dominator and from Horty's action-dominance of Remark~\ref{rem:landmarks}.)

\subsection{Sanction and goal are one machinery, toggled by sign}
\label{sec:duality}

A party's preorder does more than pick goals: its \emph{sign} on an outcome decides
whether that outcome is to be \emph{reached} or \emph{avoided}, and thereby which
necessity is in play. For a $\ppref$-\emph{good} outcome $J$ the relevant condition is
liveness, $\EF J$, and the obligations are the dominators of $J$ (``to reach the good you
must $K$''). For a $\ppref$-\emph{bad} outcome $S$---a sanction---the relevant condition
is safety, $\AG\neg S$, and the obligations are the acts without which $S$ becomes
unavoidable (``to avoid the sanction you must $K$''). Anderson's reduction of obligation
to a violation constant \src{Anderson 1958} and Kowalski and Satoh's reading of a
sanction as the worse disjunct of a goal \src{Kowalski \& Satoh 2018} already make
sanction and goal one machinery; we concede that. What we add is that a \emph{single}
party-indexed preorder, through the sign it assigns, \emph{routes the same dominator
relation} to its liveness pole or its safety pole---so Anderson's $S$ and the
teleological $G$ are not two devices but one, read at opposite polarity.

\subsection{Two orderings: Holmes and Hart on one divergence}
\label{sec:holmeshart}

The decisive move is to admit \emph{two} distinguished orderings rather than one. Let
$\plaw$ be the legal system's ordering (what one \emph{ought}, normative) and $\ppref$
the agent's own (what the agent will in fact pursue, descriptive). The same dominator
necessity, composed with each, yields two different objects:
\[
  \text{legal obligation} \;=\; \nect\circ\plaw,
  \qquad
  \text{predicted behaviour} \;=\; \nect\circ\ppref .
\]
Their relationship \emph{is} the classical jurisprudential debate. Where $\ppref$ and
$\plaw$ \emph{diverge}, the agent computes its obligations against its own ordering and
treats the law as a schedule of consequences---this is Holmes's bad man, who ``cares
only for the material consequences'' \src{Holmes 1897}. Where they \emph{coincide}, the
agent has taken the law's ordering as its own---this is Hart's internal point of view, on
which subjects ``use the rules as standards for the appraisal of \dots\ behaviour''
rather than merely predicting it \src{Hart 1961}. A sanction, on this picture, is the
law's instrument for \emph{realignment}: an attempt to depress the $\ppref$-rank of
non-compliance until compliance dominates.

We are not the first to formalise the Hart half: \src{Silk 2019} already renders Hart's
internal/external distinction as a choice of which contextual norm-set---a simplified
Kratzer premise set---a modal is read against (for Silk the contrast is endorsing vs
non-endorsing of the \emph{same} norms, not two parties' orderings); and \src{Kolt 2026}
frames precisely the Austin--Holmes (instrumental) versus Hart (internal point of view)
contrast for autonomous AI agents, tying the internal point of view to respect for the
\emph{spirit}, not merely the letter, of the law. We concede both. (By contrast,
\src{Yan \& He 2025} raise the agent-goal-vs-ideal-ordering distinction only to
\emph{collapse} it into one ideal ordering; our move is the opposite---to keep the two
apart and make their divergence the object.) Our claim is threefold: that
\emph{one} divergence does double duty as \emph{both} Holmes (where it is wide) and Hart
(where it closes); that $\ppref$ is read \emph{descriptively}, as a predictor of
behaviour, not as an endorsement set; and that the whole construction sits atop the
dominator necessity of Section~\ref{sec:nec}, so that the force-gap and the
sign-duality (below) fall out of the same schema.

\subsection{The force of an obligation as a preference gap}
\label{sec:force}

How much force does a ``must'' carry? On the present account, exactly the distance, in
the relevant ordering, between complying and breaching:
\[
  \mathrm{Force}_p(\MUST\,K)\;=\;\rank_{\ppref}(\text{comply}) - \rank_{\ppref}(\text{breach}).
\]
That force is a utility gap with a firing threshold is the economics of deterrence
\src{Becker 1968; Shavell 2003} and of qualitative decision theory \src{Pearl 1993}; we
concede the cardinal idea. Our version is \emph{ordinal} (a rank-difference over a
Kratzer preorder, not an expected utility) and \emph{party-indexed} against a paired
$\plaw$, so that the sign of the gap and the divergence of
Section~\ref{sec:holmeshart} are one phenomenon. The two pathological signs are then not
anomalies but predictions of the schema:
\begin{itemize}[noitemsep,topsep=2pt]
  \item \textbf{Gap $\le 0$ --- ``a fine is a price.''} When the sanction fails to push
  breach below compliance in $\ppref$, the ``must'' has no force; the rule is read as a
  price \src{Gneezy \& Rustichini 2000; Cooter 1984}. This corrects a tempting but
  insufficient test: force is not the \emph{magnitude} of the penalty but the
  \emph{gap} it induces, and framing or motivational crowding-out can move the gap
  without changing the magnitude.
  \item \textbf{Gap $<0$ --- inversion.} When the sanction \emph{raises} the rank of
  breach, the ``must'' inverts into a ``want to.'' Consider a protester chained to a
  railing, told ``unchain yourself or I will arrest you and you will spend the night in
  jail,'' who replies: ``you must arrest me---I came for the protest and my lodging just
  fell through.'' Jail, the intended sanction, is the protester's $\ppref$-good; the
  cancellation of his lodging is the (non-temporal) $\MAY\!\to\!\MUST$ collapse that
  makes arrest the unique remaining route to a bed. The episode also exhibits a
  Hohfeldian wrinkle the bare relation must accommodate: the act that is necessary for
  the protester's goal can be performed only by \emph{another} party, so ``you must
  arrest me'' is not an obligation he bears but a dependency on the officer's agency.
\end{itemize}

\subsection{Where the reduction is faithful}
\label{sec:faithful}

The ``bounded'' in bounded deontics is, finally, the boundary of the reduction's own
validity. Compile-time verification of the bounded \MUST{} is available exactly when
three conditions hold together: the family $\mathcal{G}$ of goals is \emph{closed}
(statically enumerable, so the property is a closed formula); the defeasible ordering on
rules is \emph{acyclic and transparent} (conflicts are resolvable where the law
genuinely prioritises---\textsc{subject to}/\textsc{notwithstanding}---yet genuine gaps
are left visible rather than silently broken); and the fact base is \emph{closed}. When
any of the three fails, the obligation falls to runtime for a correspondingly different
reason. The most important failure is the first: a goal supplied by the agent at
runtime---the anankastic ``if you want $J$''---leaves the formula open until that goal is
injected, which is also why the agent's $\ppref$, often private until revealed (as the
protester's was), is a runtime quantity whereas $\plaw$ is known in advance. Beyond these
structural limits lies a \emph{semantic} one: where the operative ordering is not the
agent's at all---Hart's internal point of view, or a categorical norm that resists the
``or else what?'' probe---the teleological reading ceases to be the faithful one, and the
reduction marks its own edge.

\subsection{Multi-party escalation (a note)}
\label{sec:atl}

Section~\ref{sec:holmeshart}'s second party invites a strategic reading: against an
adversary who may withhold the act one needs, the relevant question is not ``is $K$ on
every path?'' but ``what must $p$ do to \emph{force} $J$?''---an alternating-time
ability $\langle\!\langle p\rangle\!\rangle\,\mathsf{F}\,J$ \src{Alur, Henzinger \&
Kupferman 2002}. Obligation as such a strategic ability is \src{Wooldridge \& van der
Hoek 2005}, and our cooperative $\nect$ is the one-player (dominator) special case. We
flag this only to be honest about its status: as a standalone it is a recombination of
known parts, and we treat it as a direction for separate work rather than a claim here.
